\documentclass[11pt]{amsart}

% Encoding and fonts
\usepackage[T1]{fontenc}
\usepackage[utf8]{inputenc}
\usepackage{lmodern}

% Math packages
\usepackage{amsmath}
\usepackage{amssymb}
\usepackage{amsfonts}
\usepackage{amsthm}
\usepackage{mathtools}

% Layout and utilities
\usepackage{geometry}
\usepackage{enumitem}

\geometry{margin=1in}

\numberwithin{equation}{section}
\setcounter{tocdepth}{3}
\setcounter{secnumdepth}{3}

% ---------------------------------------------------------
% Theorem-like environments
% ---------------------------------------------------------

\theoremstyle{plain}
\newtheorem{theorem}{Theorem}[section]
\newtheorem{proposition}[theorem]{Proposition}
\newtheorem{lemma}[theorem]{Lemma}
\newtheorem{corollary}[theorem]{Corollary}
\newtheorem{conjecture}[theorem]{Conjecture}

\theoremstyle{definition}
\newtheorem{definition}[theorem]{Definition}
\newtheorem{experiment}[theorem]{Experiment}
\newtheorem{heuristic}[theorem]{Heuristic}
\newtheorem{question}[theorem]{Question}

\theoremstyle{remark}
\newtheorem{remark}[theorem]{Remark}
\newtheorem{warning}[theorem]{Warning}

\newcommand{\R}{\mathbb{R}}
\newcommand{\C}{\mathbb{C}}
\newcommand{\N}{\mathbb{N}}
\newcommand{\Hil}{\mathcal{H}}
\newcommand{\Sch}{\mathcal{S}}
\newcommand{\Four}{\mathcal{F}}
\newcommand{\Dom}{\mathcal{D}}
\newcommand{\I}{\mathrm{I}}
\newcommand{\id}{\mathrm{Id}}
\newcommand{\essran}{\operatorname{ess\,ran}}

% hyperref must be loaded last
\usepackage[hidelinks]{hyperref}

\title[Ordinal--M\"obius Quantum Operator Framework]{
Exact Operator Framework for an Ordinal--M\"obius Quantum Model:\\
M\"obius Inversion, a Zeta Operator, Generalized Zero Orbitals,\\
and Algebraic Heisenberg Non-Localization
}

\author{}
\date{}

\begin{document}

\maketitle

\tableofcontents

\bigskip

\begin{remark}[Mathematical status]
\label{rem:status}
This document is a research notebook.  The operator definitions, domain
statements, closure statements, M\"obius inversion identity, generalized
zero-mode characterization, and the algebraic non-localization theorem are
proved below.  The resulting formulation is \emph{not} a proof of the
Riemann Hypothesis.  The final Riemann-Hypothesis statement is an exact
reformulation of where the generalized null modes occur in the operator
family, and the problem of forcing all of them onto the central member of
the family remains open.
\end{remark}

\section{Introduction: what is different in the present construction}
\label{sec:introduction}

The aim is to construct a genuine operator theory in which the M\"obius
function enters arithmetically, the zeta operator is obtained by inversion,
and the non-localization statement comes from the exact canonical
commutation relation rather than from a postulated uncertainty formula.

The present construction differs from the earlier attempts in four
specific ways.

\begin{enumerate}[leftmargin=*]
\item No operator is defined by the formal expression
``$\zeta(\widehat S)$''.  The starting objects are the standard quantum
position and momentum operators and the unitary translation group generated
by momentum.

\item M\"obius enters through the norm-convergent operator series
\eqref{eq:mobius_operator_definition}.  In the half-plane where this series
is unconditionally defined in operator norm, its inverse is proved to be
the operator series \eqref{eq:zeta_operator_definition}.  The scalar zeta
Dirichlet series then appears as the Fourier multiplier of this inverse; it
is not inserted into the definition.

\item No probabilistic coarse graining of the values $-1,0,+1$ is used.
The arithmetic coefficients are exactly $\mu(n)$.

\item Heisenberg enters through the exact identity
\eqref{eq:CCR_schwartz}.  The relevant consequence is not a statement about
measurement error or variance.  It is the algebraic impossibility theorem
\ref{thm:no_joint_exact_localization}: no nonzero state, and in fact no
nonzero tempered distribution, can have simultaneously exact position and
exact momentum.
\end{enumerate}

The critical-strip operator is obtained separately from an arithmetic
alternating translation series.  Its closed realization is given with an
exact maximal domain.  Nontrivial zeros correspond to generalized null
modes of that operator.  These null modes are exact momentum modes, so
Theorem~\ref{thm:no_joint_exact_localization} makes exact spatial
localization impossible in the strict algebraic sense.

The canonical commutation relation is the standard Heisenberg relation
\cite{Heisenberg1927}.  Standard analytic facts concerning the zeta
function are used only after the arithmetic operator identities have been
established; a classical reference is \cite{Titchmarsh1986}.

\section{Quantum kinematics and exact domains}
\label{sec:quantum_kinematics}

\subsection{Hilbert space and Fourier transform}
\label{subsec:hilbert_fourier}

\begin{definition}[Hilbert space]
\label{def:hilbert_space}
Let
\begin{equation}
\Hil=L^2(\R,dq)
\label{eq:hilbert_space}
\end{equation}
with inner product
\begin{equation}
\langle \phi,\psi\rangle
=\int_{\R}\overline{\phi(q)}\psi(q)\,dq.
\label{eq:hilbert_inner_product}
\end{equation}
\end{definition}

\begin{definition}[Unitary Fourier transform]
\label{def:fourier_transform}
For $\psi\in L^1(\R)\cap L^2(\R)$ define
\begin{equation}
(\Four\psi)(p)
=\widehat\psi(p)
=\frac{1}{\sqrt{2\pi}}
\int_{\R}e^{-ipq}\psi(q)\,dq.
\label{eq:fourier_transform}
\end{equation}
The Plancherel extension is again denoted by $\Four$ and is unitary on
$\Hil$.
\end{definition}

\subsection{Position and momentum}
\label{subsec:position_momentum}

\begin{definition}[Position operator]
\label{def:position_operator_exact}
Define $Q$ by
\begin{equation}
(Q\psi)(q)=q\psi(q)
\label{eq:Q_action}
\end{equation}
with maximal domain
\begin{equation}
\Dom(Q)
=\left\{\psi\in L^2(\R):q\psi(q)\in L^2(\R)\right\}.
\label{eq:Q_domain}
\end{equation}
\end{definition}

\begin{definition}[Momentum operator]
\label{def:momentum_operator_exact}
Define $P$ by
\begin{equation}
(P\psi)(q)=-i\psi'(q)
\label{eq:P_action}
\end{equation}
with domain
\begin{equation}
\Dom(P)=H^1(\R).
\label{eq:P_domain}
\end{equation}
Equivalently,
\begin{equation}
P=\Four^{-1}M_p\Four,
\label{eq:P_fourier_realization}
\end{equation}
where $M_p$ is multiplication by $p$ with maximal multiplication domain.
\end{definition}

\begin{proposition}[Self-adjointness and common core]
\label{prop:selfadjoint_QP}
The operators $Q$ and $P$ are self-adjoint.  The Schwartz space
$\Sch(\R)$ is dense in $\Hil$, invariant under $Q$ and $P$, and is a common
operator core for both.
\end{proposition}

\begin{proof}
The operator $Q$ is the maximal multiplication operator by the real
measurable function $q$, hence is self-adjoint.  By the unitary Fourier
transform, $P$ is unitarily equivalent to the maximal multiplication
operator by the real function $p$, hence is self-adjoint.  The stated
properties of $\Sch(\R)$ follow from the defining seminorms of the Schwartz
space and standard cutoff-and-mollification approximation in the graph
norms of $Q$ and $P$.
\end{proof}

\begin{proposition}[Exact canonical commutation relation]
\label{prop:CCR}
For every $\psi\in\Sch(\R)$,
\begin{equation}
[Q,P]\psi
=(QP-PQ)\psi
=i\psi.
\label{eq:CCR_schwartz}
\end{equation}
The same identity holds on tempered distributions after extending $Q$ and
$P$ continuously from $\Sch(\R)$ to $\Sch'(\R)$ by duality.
\end{proposition}

\begin{proof}
For $\psi\in\Sch(\R)$,
\begin{align}
(QP\psi)(q)
&=-iq\psi'(q),
\label{eq:QP_computation}\\
(PQ\psi)(q)
&=-i\frac{d}{dq}\bigl(q\psi(q)\bigr)
=-i\psi(q)-iq\psi'(q).
\label{eq:PQ_computation}
\end{align}
Subtracting \eqref{eq:PQ_computation} from \eqref{eq:QP_computation} gives
\eqref{eq:CCR_schwartz}.  Polynomial multiplication and differentiation act
continuously on $\Sch(\R)$ and therefore extend by duality to
$\Sch'(\R)$; the same algebraic identity then holds distributionally.
\end{proof}

\subsection{The exact Heisenberg impossibility}
\label{subsec:exact_heisenberg_impossibility}

\begin{definition}[Exact localization]
\label{def:exact_localization}
A nonzero tempered distribution $T\in\Sch'(\R)$ is called exactly localized
at $q_0\in\R$ if
\begin{equation}
QT=q_0T.
\label{eq:exact_Q_localization}
\end{equation}
It has exact momentum $p_0\in\R$ if
\begin{equation}
PT=p_0T.
\label{eq:exact_P_localization}
\end{equation}
\end{definition}

\begin{theorem}[Algebraic Heisenberg non-localization]
\label{thm:no_joint_exact_localization}
There is no nonzero $T\in\Sch'(\R)$ satisfying both
\eqref{eq:exact_Q_localization} and \eqref{eq:exact_P_localization} for any
$q_0,p_0\in\R$.
\end{theorem}

\begin{proof}
If both eigenvalue equations held, then
\begin{align}
[Q,P]T
&=Q(p_0T)-P(q_0T)
\label{eq:joint_eigen_commutator_first}\\
&=q_0p_0T-p_0q_0T
=0.
\label{eq:joint_eigen_commutator_zero}
\end{align}
On the other hand, Proposition~\ref{prop:CCR} gives
\begin{equation}
[Q,P]T=iT.
\label{eq:joint_eigen_commutator_CCR}
\end{equation}
Equations \eqref{eq:joint_eigen_commutator_zero} and
\eqref{eq:joint_eigen_commutator_CCR} force $T=0$, contrary to the
assumption.
\end{proof}

\begin{remark}[No measurement interpretation is used]
\label{rem:no_measurement_interpretation}
Theorem~\ref{thm:no_joint_exact_localization} contains no variance, no
measurement error, and no statistical assumption.  It is a pure
non-existence statement following from the exact commutator
\eqref{eq:CCR_schwartz}.
\end{remark}

\section{Multiplicative arithmetic represented by unitary translations}
\label{sec:unitary_arithmetic}

\begin{definition}[Translation group]
\label{def:translation_group}
For $t\in\R$, define
\begin{equation}
U_t=e^{-itP}.
\label{eq:Ut_definition}
\end{equation}
Then $U_t$ is unitary on $\Hil$ and acts by
\begin{equation}
(U_t\psi)(q)=\psi(q-t)
\label{eq:Ut_action}
\end{equation}
for almost every $q$.
\end{definition}

\begin{proposition}[Fourier action and group law]
\label{prop:translation_properties}
For every $s,t\in\R$,
\begin{align}
U_tU_s
&=U_{t+s},
\label{eq:translation_group_law}\\
\Four U_t\Four^{-1}
&=M_{e^{-itp}}.
\label{eq:translation_fourier_multiplier}
\end{align}
\end{proposition}

\begin{proof}
Equation \eqref{eq:translation_group_law} follows from the functional
calculus of the self-adjoint operator $P$.  Fourier transformation of
\eqref{eq:Ut_action} gives multiplication by $e^{-itp}$, proving
\eqref{eq:translation_fourier_multiplier}.
\end{proof}

\begin{definition}[Arithmetic translations]
\label{def:arithmetic_translations}
For every $n\in\N$, define
\begin{equation}
V_n=U_{\log n}.
\label{eq:Vn_definition}
\end{equation}
\end{definition}

\begin{proposition}[Representation of the multiplicative semigroup]
\label{prop:multiplicative_representation}
The operators $V_n$ are unitary and satisfy
\begin{equation}
V_mV_n=V_{mn}
\label{eq:V_multiplicative_law}
\end{equation}
for all $m,n\in\N$.
\end{proposition}

\begin{proof}
Use \eqref{eq:translation_group_law} and
$\log m+\log n=\log(mn)$.
\end{proof}

\begin{proposition}[Ordinal--momentum commutator]
\label{prop:Q_Vn_commutator}
On $\Sch(\R)$,
\begin{equation}
[Q,V_n]=(\log n)V_n.
\label{eq:Q_Vn_commutator}
\end{equation}
\end{proposition}

\begin{proof}
Using \eqref{eq:Ut_action},
\begin{align}
(QV_n\psi)(q)
&=q\psi(q-\log n),
\label{eq:QVn_first}\\
(V_nQ\psi)(q)
&=(q-\log n)\psi(q-\log n).
\label{eq:VnQ_first}
\end{align}
Subtracting gives \eqref{eq:Q_Vn_commutator}.
\end{proof}

\section{The bounded M\"obius operator and the derived zeta operator}
\label{sec:mobius_zeta_bounded}

\subsection{Exact operator definitions for $\sigma>1$}
\label{subsec:bounded_definitions}

\begin{definition}[M\"obius operator]
\label{def:mobius_operator}
For a real parameter $\sigma>1$, define
\begin{equation}
\mathcal M_\sigma
=\sum_{n=1}^{\infty}\frac{\mu(n)}{n^\sigma}V_n.
\label{eq:mobius_operator_definition}
\end{equation}
The series is understood in the operator norm of
$\mathcal B(\Hil)$.
\end{definition}

\begin{proposition}[Norm convergence]
\label{prop:mobius_norm_convergence}
For every $\sigma>1$, the series
\eqref{eq:mobius_operator_definition} converges absolutely in operator norm
and defines a bounded operator on all of $\Hil$.  In particular,
\begin{equation}
\|\mathcal M_\sigma\|
\leq
\sum_{n=1}^{\infty}\frac{|\mu(n)|}{n^\sigma}
\leq\zeta(\sigma).
\label{eq:mobius_norm_bound}
\end{equation}
\end{proposition}

\begin{proof}
Each $V_n$ is unitary, so $\|V_n\|=1$.  The scalar majorant in
\eqref{eq:mobius_norm_bound} converges for $\sigma>1$.
\end{proof}

\begin{definition}[Candidate inverse]
\label{def:zeta_operator_bounded}
For $\sigma>1$, define
\begin{equation}
\mathcal Z_\sigma
=\sum_{n=1}^{\infty}\frac{1}{n^\sigma}V_n,
\label{eq:zeta_operator_definition}
\end{equation}
again with convergence in operator norm.
\end{definition}

\begin{theorem}[Operator M\"obius inversion]
\label{thm:operator_mobius_inversion}
For every $\sigma>1$,
\begin{equation}
\mathcal M_\sigma\mathcal Z_\sigma
=\mathcal Z_\sigma\mathcal M_\sigma
=I_{\Hil}.
\label{eq:operator_mobius_inversion}
\end{equation}
Consequently,
\begin{equation}
\mathcal Z_\sigma=\mathcal M_\sigma^{-1}.
\label{eq:zeta_as_mobius_inverse}
\end{equation}
\end{theorem}

\begin{proof}
Both series are absolutely convergent in operator norm, so their product
may be regrouped by $k=mn$.  Using
\eqref{eq:V_multiplicative_law},
\begin{align}
\mathcal M_\sigma\mathcal Z_\sigma
&=
\sum_{k=1}^{\infty}
\frac{1}{k^\sigma}
\left(\sum_{d\mid k}\mu(d)\right)V_k
\label{eq:mobius_product_regrouped}\\
&=V_1
=I_{\Hil},
\label{eq:mobius_product_identity}
\end{align}
using the M\"obius inversion identity
$\sum_{d\mid k}\mu(d)=1$ for $k=1$ and $0$ otherwise.  The same calculation
applies in the reverse order because the $V_n$ commute.
\end{proof}

\begin{remark}[The zeta operator is not defined by substitution]
\label{rem:not_zeta_of_operator}
Equation \eqref{eq:zeta_as_mobius_inverse} is the defining structural point:
$\mathcal Z_\sigma$ is the inverse of the arithmetically defined M\"obius
operator.  No expression of the form $\zeta(P)$ or $\zeta(\widehat S)$ is
used in Definition~\ref{def:zeta_operator_bounded}.
\end{remark}

\subsection{Fourier characterization}
\label{subsec:fourier_characterization}

\begin{proposition}[Exact multipliers]
\label{prop:bounded_fourier_multipliers}
For $\sigma>1$,
\begin{align}
\Four\mathcal Z_\sigma\Four^{-1}
&=M_{z_\sigma},
\qquad
z_\sigma(p)=\sum_{n=1}^{\infty}n^{-\sigma-ip},
\label{eq:zeta_multiplier_series}\\
\Four\mathcal M_\sigma\Four^{-1}
&=M_{m_\sigma},
\qquad
m_\sigma(p)=\sum_{n=1}^{\infty}\mu(n)n^{-\sigma-ip}.
\label{eq:mobius_multiplier_series}
\end{align}
The scalar series converge uniformly in $p\in\R$.  Moreover,
\begin{equation}
m_\sigma(p)z_\sigma(p)=1
\label{eq:multiplier_inverse_identity}
\end{equation}
for every real $p$.
\end{proposition}

\begin{proof}
Apply \eqref{eq:translation_fourier_multiplier} termwise; operator-norm
convergence justifies the passage to the limit.  Equation
\eqref{eq:multiplier_inverse_identity} is the Fourier transform of
\eqref{eq:operator_mobius_inversion}.
\end{proof}

\begin{definition}[Scalar zeta function in the initial half-plane]
\label{def:scalar_zeta_initial}
For $\Re s>1$, define
\begin{equation}
\zeta(s)=\sum_{n=1}^{\infty}n^{-s}.
\label{eq:scalar_zeta_dirichlet}
\end{equation}
\end{definition}

\begin{corollary}[Zeta appears as the spectral multiplier]
\label{cor:zeta_as_multiplier}
For $\sigma>1$,
\begin{align}
z_\sigma(p)
&=\zeta(\sigma+ip),
\label{eq:zeta_multiplier_identification}\\
m_\sigma(p)
&=\frac{1}{\zeta(\sigma+ip)}.
\label{eq:mobius_multiplier_identification}
\end{align}
\end{corollary}

\begin{proof}
Equation \eqref{eq:zeta_multiplier_identification} is
Definition~\ref{def:scalar_zeta_initial};
\eqref{eq:mobius_multiplier_identification} follows from
\eqref{eq:multiplier_inverse_identity}.
\end{proof}

\begin{warning}[Exact right-hand boundary]
\label{warn:sigma_gt_one_only}
The norm-convergent definitions of $\mathcal M_\sigma$ and
$\mathcal Z_\sigma$ in this section are asserted only for $\sigma>1$.
No convergence of the M\"obius operator series for $\sigma\leq1$ is assumed
or inferred.
\end{warning}

\section{The critical-strip arithmetic operator}
\label{sec:critical_strip_operator}

\subsection{Finite arithmetic operators}
\label{subsec:finite_eta_operators}

\begin{definition}[Alternating arithmetic truncations]
\label{def:eta_operator_truncations}
For $\sigma>0$ and $N\in\N$, define the bounded operator
\begin{equation}
\mathcal A_{\sigma,N}
=\sum_{n=1}^{N}(-1)^{n-1}n^{-\sigma}V_n.
\label{eq:eta_operator_truncation}
\end{equation}
Its Fourier multiplier is
\begin{equation}
a_{\sigma,N}(p)
=\sum_{n=1}^{N}(-1)^{n-1}n^{-\sigma-ip}.
\label{eq:eta_multiplier_truncation}
\end{equation}
\end{definition}

\begin{lemma}[Uniform convergence on compact momentum intervals]
\label{lem:eta_uniform_compacts}
Fix $\sigma>0$ and $R>0$.  Then $a_{\sigma,N}(p)$ converges uniformly for
$|p|\leq R$ as $N\to\infty$.
\end{lemma}

\begin{proof}
Write $s=\sigma+ip$.  Pair consecutive terms.  For $k\geq1$,
\begin{equation}
(2k-1)^{-s}-(2k)^{-s}
=s\int_{2k-1}^{2k}x^{-s-1}\,dx.
\label{eq:eta_pair_integral}
\end{equation}
For $|p|\leq R$,
\begin{equation}
\left|
(2k-1)^{-s}-(2k)^{-s}
\right|
\leq
\sqrt{\sigma^2+R^2}\,(2k-1)^{-\sigma-1}.
\label{eq:eta_pair_bound}
\end{equation}
The majorant is summable.  The possible final unpaired term has modulus at
most $N^{-\sigma}$ and tends uniformly to zero.  The Weierstrass test gives
the result.
\end{proof}

\begin{definition}[Eta multiplier]
\label{def:eta_multiplier}
For $\sigma>0$ and $p\in\R$, define
\begin{equation}
a_\sigma(p)
=\lim_{N\to\infty}a_{\sigma,N}(p)
=\sum_{n=1}^{\infty}(-1)^{n-1}n^{-\sigma-ip}.
\label{eq:eta_multiplier}
\end{equation}
Equivalently, for $\Re s>0$ define
\begin{equation}
\eta(s)=\sum_{n=1}^{\infty}(-1)^{n-1}n^{-s},
\label{eq:eta_scalar_definition}
\end{equation}
so that $a_\sigma(p)=\eta(\sigma+ip)$.
\end{definition}

\begin{remark}[Left boundary]
\label{rem:eta_left_boundary}
The direct series definition \eqref{eq:eta_scalar_definition} is used only
for $\Re s>0$.  At $\sigma=0$ the individual terms have modulus one and do
not tend to zero, so the same direct convergence statement cannot hold.
\end{remark}

\subsection{Closed realization}
\label{subsec:eta_closed_realization}

\begin{definition}[Compact-momentum core]
\label{def:compact_momentum_core}
Define
\begin{equation}
\mathcal C
=\Four^{-1}C_c^\infty(\R).
\label{eq:compact_momentum_core}
\end{equation}
Then $\mathcal C$ is dense in $\Hil$.
\end{definition}

\begin{proposition}[Strong arithmetic limit on the core]
\label{prop:eta_strong_core_limit}
For every $\sigma>0$ and every $\psi\in\mathcal C$, the limit
\begin{equation}
\mathcal A_\sigma^{(0)}\psi
=\lim_{N\to\infty}\mathcal A_{\sigma,N}\psi
\label{eq:eta_core_limit}
\end{equation}
exists in $\Hil$.  In Fourier representation,
\begin{equation}
\Four\mathcal A_\sigma^{(0)}\psi
=a_\sigma\widehat\psi.
\label{eq:eta_core_multiplier_action}
\end{equation}
\end{proposition}

\begin{proof}
The support of $\widehat\psi$ is contained in some compact interval
$[-R,R]$.  Lemma~\ref{lem:eta_uniform_compacts} gives uniform convergence
$a_{\sigma,N}\to a_\sigma$ on that interval.  Hence
\begin{equation}
\|\mathcal A_{\sigma,N}\psi-\mathcal A_\sigma^{(0)}\psi\|_2
=
\|(a_{\sigma,N}-a_\sigma)\widehat\psi\|_2
\longrightarrow0.
\label{eq:eta_core_norm_convergence}
\end{equation}
\end{proof}

\begin{definition}[Closed critical-strip operator]
\label{def:closed_eta_operator}
For $\sigma>0$, define $\mathcal A_\sigma$ to be the maximal Fourier
multiplication operator by $a_\sigma$:
\begin{equation}
\mathcal A_\sigma
=\Four^{-1}M_{a_\sigma}\Four,
\label{eq:eta_closed_operator}
\end{equation}
with exact domain
\begin{equation}
\Dom(\mathcal A_\sigma)
=\left\{
\psi\in L^2(\R):
 a_\sigma\widehat\psi\in L^2(\R)
\right\}.
\label{eq:eta_closed_domain}
\end{equation}
\end{definition}

\begin{proposition}[Closure, adjoint, and normality]
\label{prop:eta_closed_properties}
For every $\sigma>0$, $\mathcal A_\sigma$ is densely defined, closed, and
normal.  The operator $\mathcal A_\sigma^{(0)}$ of
Proposition~\ref{prop:eta_strong_core_limit} is closable and its closure is
$\mathcal A_\sigma$.  Moreover,
\begin{equation}
\mathcal A_\sigma^*
=\Four^{-1}M_{\overline{a_\sigma}}\Four
\label{eq:eta_adjoint}
\end{equation}
with
\begin{equation}
\Dom(\mathcal A_\sigma^*)
=\Dom(\mathcal A_\sigma).
\label{eq:eta_adjoint_domain}
\end{equation}
\end{proposition}

\begin{proof}
Maximal multiplication by a finite measurable function is closed and
densely defined, and its adjoint is multiplication by the complex
conjugate with the same maximal domain.  The space $C_c^\infty(\R)$ is a
core for maximal multiplication by the continuous function $a_\sigma$:
first truncate in momentum to compact sets and then approximate there by
smooth compactly supported functions.  Unitary conjugation by $\Four$
gives the stated properties.
\end{proof}

\subsection{Positive Hamiltonian}
\label{subsec:positive_hamiltonian}

\begin{definition}[Arithmetic Hamiltonian]
\label{def:arithmetic_hamiltonian}
For $\sigma>0$, define
\begin{equation}
H_\sigma
=\mathcal A_\sigma^*\mathcal A_\sigma.
\label{eq:H_sigma_definition}
\end{equation}
\end{definition}

\begin{proposition}[Exact realization of the Hamiltonian]
\label{prop:H_sigma_realization}
The operator $H_\sigma$ is positive and self-adjoint.  In Fourier
representation,
\begin{equation}
H_\sigma
=\Four^{-1}M_{|a_\sigma|^2}\Four
\label{eq:H_sigma_multiplier}
\end{equation}
with exact domain
\begin{equation}
\Dom(H_\sigma)
=\left\{
\psi\in L^2(\R):
|a_\sigma|^2\widehat\psi\in L^2(\R)
\right\}.
\label{eq:H_sigma_domain}
\end{equation}
\end{proposition}

\begin{proof}
This is the standard product of a closed densely defined operator with its
adjoint, made explicit here by the multiplication representation.  The
multiplier $|a_\sigma|^2$ is real and nonnegative, so the corresponding
maximal multiplication operator is positive self-adjoint.
\end{proof}

\section{Connection with the M\"obius-derived zeta operator}
\label{sec:eta_zeta_connection}

\begin{definition}[Dyadic factor]
\label{def:dyadic_factor}
For $\sigma>0$, define the bounded normal operator
\begin{equation}
B_\sigma
=I-2^{1-\sigma}V_2.
\label{eq:B_sigma_definition}
\end{equation}
\end{definition}

\begin{proposition}[Exact factorization for $\sigma>1$]
\label{prop:eta_zeta_factorization}
For $\sigma>1$,
\begin{equation}
\mathcal A_\sigma
=B_\sigma\mathcal Z_\sigma
=B_\sigma\mathcal M_\sigma^{-1}.
\label{eq:eta_zeta_operator_factorization}
\end{equation}
\end{proposition}

\begin{proof}
All operator series are absolutely norm-convergent in this range.  The
coefficient of $V_m$ in $B_\sigma\mathcal Z_\sigma$ is $m^{-\sigma}$ for
$m$ odd.  If $m=2n$ is even, the coefficient is
\begin{equation}
m^{-\sigma}
-2^{1-\sigma}n^{-\sigma}
=m^{-\sigma}-2m^{-\sigma}
=-m^{-\sigma}.
\label{eq:dyadic_even_coefficient}
\end{equation}
Thus the coefficient is $(-1)^{m-1}m^{-\sigma}$, which gives
\eqref{eq:eta_zeta_operator_factorization}.
\end{proof}

\begin{proposition}[Invertibility of the dyadic factor away from
$\sigma=1$]
\label{prop:B_sigma_invertibility}
If $\sigma>0$ and $\sigma\neq1$, then $B_\sigma$ is boundedly invertible and
\begin{equation}
\|B_\sigma^{-1}\|
=\frac{1}{|2^{1-\sigma}-1|}.
\label{eq:B_sigma_inverse_norm}
\end{equation}
\end{proposition}

\begin{proof}
By Fourier transformation, $B_\sigma$ is multiplication by
\begin{equation}
b_\sigma(p)
=1-2^{1-\sigma}e^{-ip\log2}.
\label{eq:B_sigma_multiplier}
\end{equation}
The minimum of $|1-re^{-i\theta}|$ over $\theta\in\R$ is $|r-1|$.  Taking
$r=2^{1-\sigma}$ proves that the multiplier is bounded away from zero
exactly when $\sigma\neq1$, and gives
\eqref{eq:B_sigma_inverse_norm}.
\end{proof}

\begin{definition}[Closed zeta operator in $0<\sigma<1$]
\label{def:critical_zeta_operator}
For $0<\sigma<1$, define
\begin{equation}
\mathfrak Z_\sigma
=B_\sigma^{-1}\mathcal A_\sigma
\label{eq:critical_zeta_operator}
\end{equation}
with
\begin{equation}
\Dom(\mathfrak Z_\sigma)
=\Dom(\mathcal A_\sigma).
\label{eq:critical_zeta_domain}
\end{equation}
\end{definition}

\begin{proposition}[Closedness and multiplier]
\label{prop:critical_zeta_properties}
For $0<\sigma<1$, $\mathfrak Z_\sigma$ is densely defined and closed.  Its
Fourier multiplier is
\begin{equation}
z_\sigma^{\mathrm{crit}}(p)
=
\frac{\eta(\sigma+ip)}
{1-2^{1-\sigma-ip}}.
\label{eq:critical_zeta_multiplier}
\end{equation}
\end{proposition}

\begin{proof}
Since $B_\sigma^{-1}$ is bounded and boundedly invertible while
$\mathcal A_\sigma$ is closed, the product
$B_\sigma^{-1}\mathcal A_\sigma$ is closed with the same domain.  The
multiplier formula follows from \eqref{eq:B_sigma_multiplier} and
\eqref{eq:eta_closed_operator}.
\end{proof}

\begin{definition}[Continuation of the scalar zeta function in the strip]
\label{def:zeta_critical_strip}
For $0<\Re s<1$, define
\begin{equation}
\zeta(s)
=\frac{\eta(s)}{1-2^{1-s}}.
\label{eq:zeta_eta_quotient}
\end{equation}
The denominator does not vanish in the open strip.
\end{definition}

\begin{proposition}[Compatibility with the initial half-plane]
\label{prop:zeta_compatibility}
For $\Re s>1$, the quotient in \eqref{eq:zeta_eta_quotient} agrees with the
Dirichlet series \eqref{eq:scalar_zeta_dirichlet}.  Consequently,
\eqref{eq:zeta_eta_quotient} gives its analytic continuation into the open
critical strip.
\end{proposition}

\begin{proof}
Absolute convergence for $\Re s>1$ allows separation of the even and odd
terms, giving
\begin{equation}
\eta(s)
=\bigl(1-2^{1-s}\bigr)\zeta(s).
\label{eq:eta_zeta_scalar_identity}
\end{equation}
The paired-series proof of Lemma~\ref{lem:eta_uniform_compacts} also gives
local uniform convergence for $\Re s>0$, hence holomorphy of $\eta$ there.
The zeros of $1-2^{1-s}$ satisfy $\Re s=1$, so the denominator is nonzero in
$0<\Re s<1$.  The identity theorem then yields the continuation.
\end{proof}

\begin{remark}[The line $\sigma=1$]
\label{rem:sigma_one_boundary}
At $\sigma=1$, $B_1$ is not boundedly invertible because its multiplier
vanishes at $p=2\pi k/\log2$.  The scalar quotient has the zeta pole at
$p=0$.  The present critical-strip construction therefore treats
$0<\sigma<1$ separately from the initial bounded region $\sigma>1$.
No bounded inverse across $\sigma=1$ is asserted.
\end{remark}

\section{Generalized zero orbitals}
\label{sec:zero_orbitals}

\subsection{Rigged Hilbert space}
\label{subsec:rigged_hilbert}

The natural setting for exact momentum modes is the standard rigged Hilbert
space
\begin{equation}
\Sch(\R)
\subset
L^2(\R)
\subset
\Sch'(\R).
\label{eq:rigged_hilbert_space}
\end{equation}

\begin{definition}[Generalized momentum orbital]
\label{def:momentum_orbital}
For $\gamma\in\R$, define the tempered distribution
\begin{equation}
\Phi_\gamma(q)
=\frac{1}{\sqrt{2\pi}}e^{i\gamma q}.
\label{eq:plane_wave_orbital}
\end{equation}
It satisfies
\begin{equation}
P\Phi_\gamma=\gamma\Phi_\gamma
\label{eq:plane_wave_momentum_eigenstate}
\end{equation}
in $\Sch'(\R)$.
\end{definition}

\begin{remark}[Not an $L^2$ vector]
\label{rem:plane_wave_not_L2}
The orbital $\Phi_\gamma$ is not in $L^2(\R)$.  It is a generalized spectral
state, exactly as an ordinary plane wave of sharp momentum in quantum
mechanics.  This distinction is essential: isolated zeros of a Fourier
multiplier do not normally produce nonzero $L^2$ kernel vectors.
\end{remark}

\subsection{Action of the arithmetic operator on exact momentum modes}
\label{subsec:eta_on_plane_waves}

\begin{proposition}[Generalized eigenvalue formula]
\label{prop:eta_plane_wave_eigenvalue}
For every $\sigma>0$ and $\gamma\in\R$, the arithmetic truncations satisfy
\begin{equation}
\mathcal A_{\sigma,N}\Phi_\gamma
=a_{\sigma,N}(\gamma)\Phi_\gamma
\label{eq:eta_truncation_plane_wave}
\end{equation}
in $\Sch'(\R)$.  Passing to the scalar limit gives the generalized action
\begin{equation}
\mathcal A_\sigma\Phi_\gamma
=\eta(\sigma+i\gamma)\Phi_\gamma.
\label{eq:eta_plane_wave_eigenvalue}
\end{equation}
\end{proposition}

\begin{proof}
From \eqref{eq:Vn_definition} and
\eqref{eq:plane_wave_orbital},
\begin{equation}
V_n\Phi_\gamma
=n^{-i\gamma}\Phi_\gamma.
\label{eq:Vn_plane_wave}
\end{equation}
Substitution into \eqref{eq:eta_operator_truncation} gives
\eqref{eq:eta_truncation_plane_wave}.  The scalar coefficients converge by
Definition~\ref{def:eta_multiplier}; multiplication of a fixed tempered
distribution by convergent scalars gives convergence in $\Sch'(\R)$.
\end{proof}

\begin{definition}[Generalized zero orbital]
\label{def:generalized_zero_orbital}
For $0<\beta<1$ and $\gamma\in\R$, the plane wave $\Phi_\gamma$ is called a
generalized zero orbital of the member $\mathcal A_\beta$ if
\begin{equation}
\mathcal A_\beta\Phi_\gamma=0
\label{eq:zero_orbital_definition}
\end{equation}
in $\Sch'(\R)$.
\end{definition}

\begin{theorem}[Zeros and generalized null modes]
\label{thm:zeros_as_generalized_null_modes}
Let $0<\beta<1$ and $\gamma\in\R$.  Then
\begin{equation}
\mathcal A_\beta\Phi_\gamma=0
\quad\Longleftrightarrow\quad
\zeta(\beta+i\gamma)=0.
\label{eq:zero_mode_equivalence}
\end{equation}
\end{theorem}

\begin{proof}
By Proposition~\ref{prop:eta_plane_wave_eigenvalue}, the left side is
equivalent to
\begin{equation}
\eta(\beta+i\gamma)=0.
\label{eq:eta_zero_equivalence_step}
\end{equation}
Inside $0<\beta<1$, the factor $1-2^{1-\beta-i\gamma}$ is nonzero.
Equation \eqref{eq:eta_zeta_scalar_identity} therefore makes
\eqref{eq:eta_zero_equivalence_step} equivalent to
$\zeta(\beta+i\gamma)=0$.
\end{proof}

\begin{corollary}[Hamiltonian form]
\label{cor:zero_orbital_hamiltonian}
If $\rho=\beta+i\gamma$ is a nontrivial zero with $0<\beta<1$, then the
generalized mode $\Phi_\gamma$ satisfies
\begin{equation}
H_\beta\Phi_\gamma=0
\label{eq:H_zero_orbital}
\end{equation}
in the generalized multiplier sense.
\end{corollary}

\begin{proof}
The Fourier multiplier of $H_\beta$ is
$|\eta(\beta+ip)|^2$, which vanishes at $p=\gamma$ by
Theorem~\ref{thm:zeros_as_generalized_null_modes}.  Multiplication of the
Dirac mass at $p=\gamma$ by this smooth local multiplier gives zero; Fourier
inversion yields \eqref{eq:H_zero_orbital}.
\end{proof}

\begin{warning}[The $L^2$ kernel is a different object]
\label{warn:L2_kernel_not_zero_set}
For a fixed $\beta$, isolated zeros of $p\mapsto\eta(\beta+ip)$ have
Lebesgue measure zero.  Therefore they do not by themselves create nonzero
vectors in $\ker_{L^2}(\mathcal A_\beta)$ or
$\ker_{L^2}(H_\beta)$.  The zero orbitals of
Definition~\ref{def:generalized_zero_orbital} are generalized spectral
modes in $\Sch'(\R)$, not normalizable bound states.
\end{warning}

\section{Where exact Heisenberg non-localization enters}
\label{sec:heisenberg_zero_orbitals}

\begin{theorem}[A zero orbital cannot be exactly localized]
\label{thm:zero_orbital_nonlocalization}
Let $\rho=\beta+i\gamma$ be a nontrivial zero with $0<\beta<1$.  Its
generalized zero orbital $\Phi_\gamma$ cannot satisfy
\begin{equation}
Q\Phi_\gamma=q_0\Phi_\gamma
\label{eq:forbidden_zero_localization}
\end{equation}
for any $q_0\in\R$.
\end{theorem}

\begin{proof}
By Definition~\ref{def:momentum_orbital},
$P\Phi_\gamma=\gamma\Phi_\gamma$.  If
\eqref{eq:forbidden_zero_localization} also held, then
$\Phi_\gamma$ would be a nonzero simultaneous generalized eigenvector of
$Q$ and $P$, contradicting
Theorem~\ref{thm:no_joint_exact_localization}.
\end{proof}

\begin{corollary}[Exact, not statistical, non-localization]
\label{cor:exact_not_statistical}
For every nontrivial zero $\rho=\beta+i\gamma$, the associated exact
momentum orbital is algebraically forbidden from being an exact position
orbital.  The obstruction is the identity $[Q,P]=iI$ and does not depend on
an uncertainty variance or on a measurement model.
\end{corollary}

\begin{remark}[What is and is not quantized]
\label{rem:beta_parameter_gamma_momentum}
In the present rigorous model, $\gamma$ is the generalized eigenvalue of
the self-adjoint momentum $P$, while $\beta$ labels the member
$\mathcal A_\beta$ of the arithmetic operator family.  The theory does
\emph{not} yet promote $\beta$ to a second self-adjoint observable.  This is
deliberate: doing so without defining a domain and a compatible canonical
structure would reintroduce a merely formal ``hat on a variable''.
\end{remark}

\section{Operator form of the Riemann Hypothesis}
\label{sec:operator_RH}

\begin{theorem}[Exact reformulation of RH]
\label{thm:operator_RH_equivalence}
The Riemann Hypothesis is equivalent to the following operator statement:
for every $0<\sigma<1$ and every $\gamma\in\R$,
\begin{equation}
\mathcal A_\sigma\Phi_\gamma=0
\quad\Longrightarrow\quad
\sigma=\frac12.
\label{eq:operator_RH_statement}
\end{equation}
Equivalently, all generalized momentum-null modes of the critical-strip
family occur in the central member $\mathcal A_{1/2}$.
\end{theorem}

\begin{proof}
By Theorem~\ref{thm:zeros_as_generalized_null_modes},
$\mathcal A_\sigma\Phi_\gamma=0$ is equivalent to
$\zeta(\sigma+i\gamma)=0$ in the open critical strip.  Therefore
\eqref{eq:operator_RH_statement} is exactly the assertion that every
nontrivial zero has real part $1/2$.
\end{proof}

\begin{warning}[No proof of RH has occurred]
\label{warn:no_RH_proof}
Theorem~\ref{thm:operator_RH_equivalence} is an equivalence, not a proof of
\eqref{eq:operator_RH_statement}.  The exact Heisenberg theorem proves that
a zero orbital cannot be simultaneously localized in $Q$ and $P$; by
itself it does not force the family parameter $\sigma$ to equal $1/2$.
Any argument claiming otherwise would require an additional theorem tying
$\sigma$ to the canonical pair.
\end{warning}

\section{Finite operators for exact numerical work}
\label{sec:finite_numerics}

\begin{definition}[Finite M\"obius and alternating operators]
\label{def:finite_operators}
For $N\in\N$ and real $\sigma$, define bounded operators on all of $\Hil$ by
\begin{align}
\mathcal M_{\sigma,N}
&=\sum_{n=1}^{N}\mu(n)n^{-\sigma}V_n,
\label{eq:finite_mobius_operator}\\
\mathcal A_{\sigma,N}
&=\sum_{n=1}^{N}(-1)^{n-1}n^{-\sigma}V_n.
\label{eq:finite_alternating_operator}
\end{align}
\end{definition}

\begin{proposition}[Finite Fourier multipliers]
\label{prop:finite_multiplier_form}
The operators in Definition~\ref{def:finite_operators} are bounded normal
operators with Fourier multipliers
\begin{align}
m_{\sigma,N}(p)
&=\sum_{n=1}^{N}\mu(n)n^{-\sigma-ip},
\label{eq:finite_mobius_multiplier}\\
a_{\sigma,N}(p)
&=\sum_{n=1}^{N}(-1)^{n-1}n^{-\sigma-ip}.
\label{eq:finite_eta_multiplier}
\end{align}
\end{proposition}

\begin{proof}
Each is a finite linear combination of commuting unitary Fourier
multipliers $V_n$.  Equations \eqref{eq:finite_mobius_multiplier} and
\eqref{eq:finite_eta_multiplier} follow from
\eqref{eq:translation_fourier_multiplier}.
\end{proof}

\begin{remark}[Purpose of finite truncations]
\label{rem:finite_truncations_purpose}
The finite operators contain no domain ambiguity and can be used to test
spectral profiles numerically.  Such experiments do not replace the closed
operator definitions of Sections~\ref{sec:mobius_zeta_bounded} and
\ref{sec:critical_strip_operator}.
\end{remark}

\section{What has been established and what remains open}
\label{sec:open_questions}

\subsection{Rigorous statements established in this notebook}
\label{subsec:rigorous_summary}

The following statements have been proved within the present framework.

\begin{enumerate}[leftmargin=*]
\item $Q$ and $P$ are self-adjoint on the exact domains
\eqref{eq:Q_domain} and \eqref{eq:P_domain}, with the exact canonical
commutator \eqref{eq:CCR_schwartz} on $\Sch(\R)$.

\item No nonzero tempered distribution can possess simultaneously exact
$Q$ and exact $P$; this is Theorem~\ref{thm:no_joint_exact_localization}.

\item The arithmetic translations $V_n$ form a unitary representation of
the multiplicative semigroup of positive integers.

\item For $\sigma>1$, the M\"obius operator
$\mathcal M_\sigma$ is bounded and invertible, and its inverse is the zeta
operator $\mathcal Z_\sigma$ by the exact operator M\"obius inversion
identity \eqref{eq:operator_mobius_inversion}.

\item For every $\sigma>0$, the alternating arithmetic operator
$\mathcal A_\sigma$ has the closed maximal realization
\eqref{eq:eta_closed_operator}--\eqref{eq:eta_closed_domain}.

\item In the open critical strip, nontrivial zeros are exactly generalized
null momentum modes of $\mathcal A_\sigma$.

\item Every such exact momentum mode is algebraically forbidden by
Heisenberg from being an exact position mode.
\end{enumerate}

\subsection{Open problem that would be needed for an RH mechanism}
\label{subsec:RH_mechanism_open}

\begin{question}[Central structural question]
\label{q:central_structural_question}
Is there a property derivable from the M\"obius-generated operator structure
itself that distinguishes $\mathcal A_{1/2}$ from all
$\mathcal A_\sigma$, $\sigma\neq1/2$, and forces every generalized null
momentum mode of the family to occur at $\sigma=1/2$?
\end{question}

\begin{remark}
\label{rem:heisenberg_not_enough_for_RH}
The exact Heisenberg obstruction is necessary for the intended
non-localization interpretation, but it is not sufficient to answer
Question~\ref{q:central_structural_question}.  A successful next step must
produce a new structural link between the family parameter $\sigma$ and
the canonical pair $(Q,P)$; it cannot be obtained merely by rewriting
$\zeta(s)$ in operator notation.
\end{remark}

\begin{question}[Possible strengthening]
\label{q:beta_operator_question}
Can the real coordinate $\sigma-1/2$ be promoted to a genuinely defined
self-adjoint observable on an enlarged Hilbert space, with a specified
domain and an exact commutation law, while preserving the M\"obius inversion
construction of $\mathcal Z_\sigma$?  Any such construction must be given
before assigning a quantum interpretation to the real part of a zero.
\end{question}

\section*{Notebook History}

\noindent Version 0.1: Exact Hilbert-space domains for $Q$ and $P$, the
unitary arithmetic translations $V_n$, bounded M\"obius inversion for
$\sigma>1$, the closed critical-strip alternating operator for $\sigma>0$,
generalized zero orbitals, and the algebraic Heisenberg non-localization
theorem assembled into a single framework.

\begin{thebibliography}{99}

\bibitem{Heisenberg1927}
W.~Heisenberg,
\emph{"Uber den anschaulichen Inhalt der quantentheoretischen Kinematik und
Mechanik},
Zeitschrift f\"ur Physik \textbf{43} (1927), 172--198.
doi:10.1007/BF01397280.

\bibitem{Titchmarsh1986}
E.~C.~Titchmarsh,
\emph{The Theory of the Riemann Zeta-Function},
2nd ed., revised by D.~R.~Heath-Brown,
Oxford University Press, Oxford, 1986.

\end{thebibliography}

\end{document}
